\documentclass[12pt]{article}%
\usepackage{amsfonts}
\usepackage{fancyhdr}
\usepackage{comment}
\usepackage[a4paper, top=2.5cm, bottom=2.5cm, left=2.2cm, right=2.2cm]%
{geometry}
\usepackage{times}
\usepackage{amsmath}
\usepackage{changepage}
\usepackage{amssymb}
\usepackage{graphicx}%
\usepackage{hyperref}
\setcounter{MaxMatrixCols}{30}
\newtheorem{theorem}{Theorem}
\newtheorem{acknowledgement}[theorem]{Acknowledgement}
\newtheorem{algorithm}[theorem]{Algorithm}
\newtheorem{axiom}{Axiom}
\newtheorem{case}[theorem]{Case}
\newtheorem{claim}[theorem]{Claim}
\newtheorem{conclusion}[theorem]{Conclusion}
\newtheorem{condition}[theorem]{Condition}
\newtheorem{conjecture}[theorem]{Conjecture}
\newtheorem{corollary}[theorem]{Corollary}
\newtheorem{criterion}[theorem]{Criterion}
\newtheorem{definition}[theorem]{Definition}
\newtheorem{example}[theorem]{Example}
\newtheorem{exercise}[theorem]{Exercise}
\newtheorem{lemma}[theorem]{Lemma}
\newtheorem{notation}[theorem]{Notation}
\newtheorem{problem}[theorem]{Problem}
\newtheorem{proposition}[theorem]{Proposition}
\newtheorem{remark}[theorem]{Remark}
\newtheorem{solution}[theorem]{Solution}
\newtheorem{summary}[theorem]{Summary}
\newenvironment{proof}[1][Proof]{\textbf{#1.} }{\ \rule{0.5em}{0.5em}}

\newcommand{\Q}{\mathbb{Q}}
\newcommand{\R}{\mathbb{R}}
\newcommand{\C}{\mathbb{C}}
\newcommand{\Z}{\mathbb{Z}}

\begin{document}

\title{CS690S -AI Alignment - Fall 2026}
\author{Sagnik Chatterjee}
\date{\today}
\maketitle


\section{Part 1}

\subsection{Statement}
Run the following code again
\begin{verbatim}
python test_gym.py
\end{verbatim}
What do you notice happening? What specifically about MountainCar makes the
problem difficult for RL?

\subsection{Solution}
	YOUR SOLUTION GOES HERE




\section{Part 2}

\subsection{Statement}
You will now learn how to solve the MountainCar task by driving the car yourself.
Run
\begin{verbatim}
python mountain_car_play.py
\end{verbatim}
Use the arrow keys to control the car by accelerating left and right. Note that
if you run out of time you get a reward of $-200.0$ and the car will reset to the
bottom of the hill. If you get to the flag in less than 200 steps you can get a
higher score. Once you reach the flag the environment will restart at the bottom
of the hill.

Keep practicing until you can reliably get out of the valley. You can see your
score for each episode output on the command line.
Experiment with different strategies: e.g.
\begin{enumerate}
	\item Going right, then left, then right all the way up the hill
	\item Going left, then right all the way up the hill.
	\item Left, right, left, right up the hill.
\end{enumerate}
Which strategy do you like best? What is the best score you can get as a human
demonstrator?

\subsection{Solution}
	YOUR SOLUTION GOES HERE




\section{Part 3}

\subsection{Statement}
Now we will teach a simple behavioral cloning (BC) agent to drive itself out of
the valley.

First take a look at \texttt{mountain\_car\_bc.py} and try to get a basic
understanding of what is happening. By default this code will collect a single
demonstration, then parse the states and actions into tensors for easy input to a
PyTorch neural network. It then trains a policy to imitate the demonstrations and
evaluates the policy by testing it on different initial states. The command line
will output the average, min, and max total rewards.

Try it out by running
\begin{verbatim}
python mountain_car_bc.py
\end{verbatim}
Try to provide a good demonstration. Then watch what the agent has learned. Does
it do a good job imitating? Does it ever get stuck? What is the average, min, and
max?

\subsection{Solution}
	YOUR SOLUTION GOES HERE




\section{Part 4}

\subsection{Statement}
Let's give more than one demonstration. Run the following to give 5 good
demonstrations. If you mess up during one demo, feel free to restart until you
give 5 good demos and try to keep to a consistent strategy.
\begin{verbatim}
python mountain_car_bc.py --num_demos 5
\end{verbatim}
Report the average min and max returns. Did it perform any better? Why or why
not? Does the agent copy the strategy you used?

\subsection{Solution}
	YOUR SOLUTION GOES HERE




\section{Part 5}

\subsection{Statement}
What do you think will happen if we give good and bad demonstrations?
You will now give two demonstrations. For the first one, just press the right
arrow key for the entire episode until it restarts. Then for the second demo,
give a good demonstration that quickly gets out of the valley.
\begin{verbatim}
python mountain_car_bc.py --num_demos 2
\end{verbatim}
What does the policy learn? Why might bad demonstrations be a problem? Briefly
suggest one potential idea for making BC robust to bad demonstrations, as long as
they are a minority of demonstrations.

\subsection{Solution}
	YOUR SOLUTION GOES HERE




\section{Part 6}

\subsection{Statement}
At a conceptual level, describe what changes you would need to make to the BC
code to implement BCO(0), as described in the paper we read in class:
\url{https://arxiv.org/pdf/1805.01954.pdf}. Answer this question before starting
Part 7.

\subsection{Solution}
	YOUR SOLUTION GOES HERE




\section{Part 7}

\subsection{Statement}
Implement and test BCO(0) by modifying the starter code in
\texttt{mountain\_car\_bco.py}. This starter code is identical to the BC code,
except: (1) the function \texttt{collect\_random\_interaction\_data} has been
added to help with learning the inverse dynamics model; (2) a new argument for
number of inverse dynamics training iterations has been added; and (3) some
comments have been added about where to use the inverse dynamics model instead of
ground-truth actions. You may need to experiment with different neural network
sizes and training iterations to get the inverse dynamics model to work well.
Report how well it works and what you tried to get it to work, including stats on
the accuracy of prediction when tested on the demonstration data (training should
be done with the random interaction data, not the demonstrations).

\subsection{Solution}
	YOUR SOLUTION GOES HERE




\end{document}
